\documentclass[aspectratio=169]{beamer}  % 16:9 aspect ratio

% Use a clean theme as base
\usetheme{default}
\usecolortheme{default}
\useoutertheme{miniframes}
\useinnertheme{circles}
\usepackage{booktabs}
\usepackage[backend=bibtex,sorting=none]{biblatex}
\addbibresource{ref.bib} 
\setbeamerfont{footnote}{size=\tiny}

% Custom colors from HKUST logo
\definecolor{hkustblue}{RGB}{0, 51, 119}    % Navy blue from logo
\definecolor{hkustgold}{RGB}{180, 141, 61}  % Golden brown from logo
\definecolor{lightgray}{RGB}{236, 240, 241}


% Customize the appearance
\setbeamercolor{structure}{fg=hkustblue}
\setbeamercolor{background canvas}{bg=white}
\setbeamercolor{normal text}{fg=hkustblue}
\setbeamercolor{frametitle}{fg=hkustblue,bg=white}
\setbeamercolor{itemize item}{fg=hkustgold}
\setbeamercolor{itemize subitem}{fg=hkustgold}
\setbeamercolor{block title}{fg=white,bg=hkustblue}
\setbeamercolor{block body}{fg=hkustblue,bg=lightgray}
\setbeamercolor{title}{fg=hkustblue}
\setbeamercolor{author}{fg=hkustgold}


\setbeamertemplate{title page}
{
  \begin{centering}
    \vspace{1cm}
    {\usebeamercolor[fg]{title}\LARGE\inserttitle\par}
    \vspace{0.8cm}
    {\usebeamercolor[fg]{author}\large\insertauthor\par}
    \vspace{0.8cm}
    {\usebeamercolor[fg]{institute}\normalsize\insertinstitute\par}
    \vspace{0.2cm}
    {\usebeamercolor[fg]{date}\normalsize\insertdate\par}
    \vfill
  \end{centering}
}

% Custom command to show navigation symbols after title and outline pages
\makeatletter
\def\shownavigationsymbols{%
    \ifnum\c@framenumber>2%
        \setbeamertemplate{navigation symbols}[default]%
    \else%
        \setbeamertemplate{navigation symbols}{}%
    \fi%
}
\addtobeamertemplate{footline}{\shownavigationsymbols}{}
\makeatother

% Customize frame title
\setbeamertemplate{frametitle}{
    \insertframetitle
}

% Customize itemize bullets
\setbeamertemplate{itemize item}{\small\raise0.5pt\hbox{\textbullet}}
\setbeamertemplate{itemize subitem}{\tiny\raise1.5pt\hbox{\textbullet}}

% Packages
\usepackage{graphicx}
\usepackage{amsmath}
\usepackage{hyperref}
\usepackage{multirow}
\usepackage{array}
\usepackage{comment}

% Paper details (replace with your information)
\title{LCoffee}
%\author{Kamenica \& Gentzkow, AER, 2011}
%\institute{Presented by \\ Tao LI}
\date{\today}


% Mark the page number (without any other info) in the footer except the title page
\setbeamertemplate{footline}{%
    \ifnum\value{framenumber}>1
    \hfill\insertframenumber\fi
}


\begin{document}

% Title page
\begin{frame}
    \titlepage
\end{frame}

% Table of contents
\begin{frame}{Outline}
    \tableofcontents
\end{frame}

% Get a intro page of each section
% \AtBeginSection[]
% {
%     \begin{frame}
%         \frametitle{Outline}
%         \tableofcontents[currentsection]
%     \end{frame}
% }

% Mark the page number on a specific position
\addtobeamertemplate{navigation symbols}{}{%
    \usebeamerfont{footline}%
    \usebeamercolor[fg]{footline}%
    \hspace{1em}%
    \insertframenumber/\inserttotalframenumber
}


\section{Updates}
\begin{frame}{Two Outcome Variables}
\small
The four update regressions use two families of outcome variables $Y_{iet}$:

\begin{table}[h]
\centering
\scriptsize
\begin{tabular}{p{0.22\linewidth} p{0.64\linewidth}}
\toprule
\textbf{Outcome type} & \textbf{Definition} \\
\midrule
Purchase frequency &
$n\_purchases =$ purchases in the event-time window divided by window length (daily rate) \\
Novelty seeking &
share of purchases that are ``new'' (never purchased before) within the window \\
\bottomrule
\end{tabular}
\end{table}

\begin{itemize}
    \item Purchase-frequency slides ask whether closures change \emph{how much} members buy afterward
    \item Novelty-seeking slides ask whether closures change \emph{what kinds of products} members explore afterward
\end{itemize}
\end{frame}



\begin{frame}{Update 1: Matched Purchase DiD}
Matched panel of members observed in every pre/post period, with
$Y_{iet}=$ daily purchase rate and $t \in \{-4,-3,-2,-1,1,2,3,4\}$:
\begingroup
\small
\[
Y_{iet} =
\delta \cdot \mathbf{1}(\text{post}_{t}) \times \mathbf{1}(\text{Treated}_{ie})
+ \phi_{ie} + \omega_t + \gamma_m + \varepsilon_{iet}
\]
\endgroup

\begin{columns}[T]
\begin{column}{0.41\linewidth}
\textbf{Parameter of interest}
\[
\hat{\delta} = 0.00514 \; (0.00441)
\]
\end{column}
\begin{column}{0.56\linewidth}
\textbf{Result}
\begin{itemize}
    \item $p=0.244$: insignificant
    \item Treated / control members: 952 / 4,248
    \item 22 closures, 5,200 members, 41,600 obs.
    \item Pre-trend joint test: $p=0.458$
\end{itemize}
\end{column}
\end{columns}
\end{frame}

\begin{frame}{Update 2: Purchase DDD}
Unmatched-panel DDD for purchase frequency, where $D_{ie}$ is the displacement indicator:
\begingroup
\small
\[
\begin{aligned}
Y_{iet} &= \delta^B \bigl(\text{post}_{t} \times \text{Treated}_{ie}\bigr)
+ \beta \bigl(\text{post}_{t} \times D_{ie}\bigr) \\
&\quad + \delta^D \bigl(\text{post}_{t} \times \text{Treated}_{ie} \times D_{ie}\bigr)
+ \phi_{ie} + \omega_t + \gamma_m + \varepsilon_{iet}
\end{aligned}
\]
\endgroup

\begin{columns}[T]
\begin{column}{0.42\linewidth}
\textbf{Estimated parameters}
\begin{table}[h]
\centering
\scriptsize
\begin{tabular}{lc}
\toprule
Parameter & Estimate (SE) \\
\midrule
$\hat{\delta}^B$ & 0.00081 (0.00064) \\
$\hat{\beta}$ & $-0.03353^{***}$ (0.00107) \\
$\hat{\delta}^D$ & $0.00517^{**}$ (0.00241) \\
\bottomrule
\end{tabular}
\end{table}
\end{column}
\begin{column}{0.57\linewidth}
\textbf{Result}
\begin{itemize}
    \item Treated / control members: 8,092 / 36,766
    \item Baseline effect is near zero
    \item Displacement effect is small and positive
    \item Pre-trend didn't pass $p=0.004$
\end{itemize}
\end{column}
\end{columns}
\end{frame}

\begin{frame}{Update 3: Matched Novelty-Seeking DiD}
Matched-panel DiD for novelty seeking:
\begingroup
\small
\[
Y_{iet} =
\delta \cdot \mathbf{1}(\text{post}_{t}) \times \mathbf{1}(\text{Treated}_{ie})
+ \phi_{ie} + \omega_t + \gamma_m + \varepsilon_{iet}
\]
\endgroup

\begin{columns}[T]
\begin{column}{0.45\linewidth}
\textbf{Estimated treatment effect}
\begin{table}[h]
\centering
\scriptsize
\begin{tabular}{lc}
\toprule
Outcome mode & $\hat{\delta}$ \\
\midrule
Distinct products & $0.04066^{***}$ (0.01419) \\
Purchase instances & $0.02991^{**}$ (0.01435) \\
\bottomrule
\end{tabular}
\end{table}
\end{column}
\begin{column}{0.53\linewidth}
\textbf{Result}
\begin{itemize}
    \item Treated / control members: 308 / 1,415
    \item Positive post-closure effect in both outcome constructions
    \item Pre-trends supported: $p=0.525$ and $0.686$
\end{itemize}
\end{column}
\end{columns}
\end{frame}

\begin{frame}{Update 4: Novelty-Seeking DDD}
Unmatched-panel DDD for distinct-mode novelty seeking:
\begingroup
\small
\[
\begin{aligned}
Y_{iet} &= \delta^B \bigl(\text{post}_{t} \times \text{Treated}_{ie}\bigr)
+ \beta \bigl(\text{post}_{t} \times D_{ie}\bigr) \\
&\quad + \delta^D \bigl(\text{post}_{t} \times \text{Treated}_{ie} \times D_{ie}\bigr)
+ \phi_{ie} + \omega_t + \gamma_m + \varepsilon_{iet}
\end{aligned}
\]
\endgroup

\begin{columns}[T]
\begin{column}{0.4\linewidth}
\textbf{Estimated parameters}
\begin{table}[h]
\centering
\scriptsize
\begin{tabular}{lc}
\toprule
Parameter & Estimate (SE) \\
\midrule
$\hat{\delta}^B$ & $0.031^{***}$ (0.011) \\
$\hat{\beta}$ & $0.036^{***}$ (0.006) \\
$\hat{\delta}^D$ & $-0.039^{***}$ (0.014) \\
\bottomrule
\end{tabular}
\end{table}
\end{column}
\begin{column}{0.57\linewidth}
\textbf{Result}
\begin{itemize}
    \item Treated / control members: 8,092 / 36,766
    \item Baseline effect is positive for non-displaced treated members
    \item Displacement effect is negative
\end{itemize}
\end{column}
\end{columns}
\end{frame}


\begin{frame}{Results}
\textbf{The Core Narrative}: Forced non-purchase activates an intense, narrowly-targeted goal-completion drive. This drive makes displaced customers buy more and explore less when the store reopens.
\small
\begin{itemize}
    \item The forced not to purchase effect on whose who want to purchase makes them more likely to buy afterwards.
    \begin{itemize}
        \item The large negative \texttt{post x disp} term indicates the existence of a purchase circle.
        \item In the novelty-seeking DDD, $\delta^d$ is positive.
    \end{itemize}    
    \item The forced not to purchase effect on whose who want to purchase makes them less novelty-seeking afterwards.
    \begin{itemize}
        \item In the novelty-seeking DDD, $\delta^d$ is significantly negative.
    \end{itemize}
\end{itemize}
\end{frame}

% ============================================================
%  Frame 1: Conceptualizing "Displacement"
% ============================================================
\begin{frame}{Conceptualizing Displacement: Active Purchase Goal}
\footnotesize

\textbf{Core Idea.}
The displacement indicator $D_{ie}$ is a \emph{revealed-preference proxy} for an
\textbf{active purchase goal}---a chronically activated consumption intention
made visible by customers' purchasing pattern.

\vspace{0.5cm}

\medskip
\tiny
\begin{tabular}{@{} p{0.30\linewidth} p{0.64\linewidth} @{}}
\toprule
\textbf{Concept} & \textbf{Key Statement \& Citation} \\
\midrule
Habitual goal pursuit &
Habitual behavior is fundamentally goal-driven; recurring
purchase routines reflect a chronically active consumption goal,
not mere inertia.
\newline
\textit{Kruglanski \& Szumowska (2020, Perspectives on Psychological Science)} \\
\addlinespace[4pt]
 &
Goals are mental representations that, once activated,
remain accessible and direct behavior through a network
of means--ends associations.
\newline
\textit{Kruglanski et al.\ (2002, Advances in Experimental Social Psychology)} \\
\addlinespace[4pt]
Implemental mindset &
An initial purchase shifts the consumer from a deliberative
to an implemental mindset, making subsequent purchases more
likely---consecutively purchasing customers are already in
this action-oriented state.
\newline
\textit{Dhar, Huber \& Khan (2007, Journal of Marketing Research);
Gollwitzer (1990)} \\
\bottomrule
\end{tabular}
\end{frame}


% ============================================================
%  Frame 2: Why Closure Affects Displaced Customers
% ============================================================
\begin{frame}{Conceptualize Closure Impact on Displaced Customers}
\footnotesize

\textbf{Core Idea.}
When a store closure blocks an active purchase goal, the goal
persists cognitively and intensifies, producing a narrowly-targeted
rebound upon reopening.

\vspace{0.5cm}

\medskip
\tiny
\begin{tabular}{@{} p{0.26\linewidth} p{0.68\linewidth} @{}}
\toprule
\textbf{Mechanism} & \textbf{Key Statement \& Citation} \\
\midrule
Goal persistence (Zeigarnik effect) &
Unfulfilled goals remain mentally active---causing intrusive
thoughts and heightened accessibility---until they are completed
or a concrete plan is formed.
\newline
\textit{Masicampo \& Baumeister (2011, JPSP)} \\
\addlinespace[4pt]
Mental simulation \& counterfactual thinking &
Blocked goals trigger counterfactual and prefactual thinking
(``I would be buying\ldots''), which is closely connected to
goal cognition and regulates subsequent behavior via both
content-specific and motivational pathways.
\newline
\textit{Epstude \& Roese (2008, PSPR);
Kappes (2016, SPPC)} \\
\addlinespace[4pt]
Psychological reactance &
Restricting consumers' freedom to choose triggers reactance---a
motivational state that intensifies desire for the restricted
option and drives efforts to restore the threatened freedom.
\newline
\textit{Brehm (1966); Clee \& Wicklund (1980, JCR)} \\
\addlinespace[4pt]
Consumer response to unavailability &
As personal commitment to an unavailable option increases,
consumers react more negatively and are more motivated
to fulfill the thwarted purchase.
\newline
\textit{Fitzsimons (2000, JCR)} \\
\bottomrule
\end{tabular}
\end{frame}

\section{Mechanism}
\begin{frame}{Mechanism 1: Goal Carryover}
\footnotesize
\textbf{Idea}
\begin{itemize}
    \item Some customers intended to buy during the closure, but were temporarily blocked.
    \item After reopening, these unsatisfied purchase goals generate a short-run rebound in purchase prob.
\end{itemize}

\textbf{Supportive evidence 1}: Positive and significant immediately after reopening

\begin{table}
\scriptsize
\centering
\begin{tabular}{lcc}
\toprule
Period & $\delta_\ell$ & $p$ \\
\midrule
$-4$ & 0.00043 & 0.720 \\
$-3$ & -0.00028 & 0.809 \\
$-2$ & -0.00031 & 0.759 \\
$1$  & 0.00254$^{**}$ & 0.014 \\
$2$  & 0.00016 & 0.884 \\
$3$  & 0.00623$^{***}$ & $<0.001$ \\
$4$  & 0.00175 & 0.130 \\
\bottomrule
\end{tabular}
\end{table}

\textbf{Supportive evidence 2}
\begin{itemize}
    \item 7.3\% of selected treatment members purchased at another store during the closure.
\end{itemize}
\end{frame}

\begin{frame}{Mechanism 2: Missed Exposure to New Products}
\footnotesize
\textbf{Idea}
\begin{itemize}
    \item Control customers keep purchasing during the closure, so they continue to see and try newly introduced products.
    \item Treated high-intention customers are interrupted and miss this exposure window.
    \item After reopening, treated customers may be less familiar with the updated assortment and return to safer, previously chosen products.
    \item This can lower post-closure novelty-seeking for displaced treated customers relative to similarly high-intention controls.
\end{itemize}

\textbf{Interpretation}
\begin{itemize}
    \item The negative novelty-seeking DDD effect may partly reflect missed exposure to assortment changes during the closure period.
\end{itemize}
\end{frame}

\begin{frame}{Mechanism 2: Potential Tests}
\footnotesize
\textbf{Descriptive evidence}
\begin{itemize}
    \item Across 22 closures, the matched control stores introduced on average 7.9 new products per store during the closure.
    \item The median average is 8.6, with a range from 1.0 to 15.4.
\end{itemize}

\textbf{Potential tests}
\begin{itemize}
    \item Test whether the negative novelty-seeking displacement effect is stronger when matched control stores experienced more new-product introductions during the closure.
    \item Test whether treated customers who purchased at other stores during the closure show a weaker decline in novelty-seeking.
    \item Use product introduction timing to test whether treated customers are less likely to adopt products first introduced during the closure period.
\end{itemize}
\end{frame}

\begin{comment}

% ------ Paper 1: Soysal et al. (2019) ------
\begin{frame}{\textbf{Store Closure $\rightarrow$ Customer Churn} (\textit{POM, 2019})}
\footnotesize
Gonca Soysal, Alejandro Zentner, and Zhiqiang (Eric) Zheng
\vspace{4pt}

Setting \& Data
\begin{itemize}
    \item DVD rental company offering both online-by-mail subscriptions and physical store locations
    \item Individual-level panel: subscription status, online/offline rental transactions
    \item Large wave of physical store closures during the study period
\end{itemize}
 
Identification
\begin{itemize}
    \item Store closure treated as exogenous
    \item Logistic regression
\end{itemize}

Results
\begin{itemize}
    \item Store closure in a zip code increased the churn rate by $\approx 50\%$ for \textbf{heavy} users of physical stores
    \item Increase in churn was essentially \textbf{zero} for light users
    \item Physical stores substantially increase loyalty, but only for a fraction of consumers $\Rightarrow$ substantial heterogeneity in how much utility consumers derive from having the physical-store option
\end{itemize}

\end{frame}

 
% ------ Paper 5: Ergin, Gümüş & Yang (2022) ------
\begin{frame}{\textbf{Intra-Firm Product Substitutability} (\textit{POM, 2022})}
\footnotesize
Yang El\c{c}in Ergin, Mehmet G\"{u}m\"{u}\c{s}, Nathan Yang
\vspace{4pt}
 
Setting \& Data
\begin{itemize}
    \item Large fast-fashion retailer with 300+ stores
    \item Weekly sales and inventory levels for all products across the entire store network
\end{itemize}
 
Identification
\begin{itemize}
    \item Instrumented Difference-in-Differences (DDIV): instruments for stock-outs to address endogeneity
    \item Exploits within-firm, cross-store variation in product availability
\end{itemize}
 
Results
\begin{itemize}
    \item Sales of an item at a focal store \textbf{increase} when that item stocks out at neighboring stores
    \item Substitutability \textbf{strongest} at first stock-out observation; \textbf{decays} over time
    \item More pronounced for \textbf{street} stores vs.\ mall stores
\end{itemize}

\end{frame}
 
% ------ Paper 2: Tohidi et al. (2022) ------
\begin{frame}{\textbf{Habits in Consumer Purchases} (\textit{WP, 2022})}
\footnotesize
Amir Tohidi, Dean Eckles, Ali Jadbabaie
\vspace{4pt}
 
Setting \& Data
\begin{itemize}
    \item NielsenIQ household panel — purchases across many CPG categories at U.S.\ grocery stores
    \item Store closures as exogenous shocks disrupting the familiar shopping context
\end{itemize}
 
Identification
\begin{itemize}
    \item DiD: treated = households exposed to a closure, with four exposure groups (E1--E4) based on purchase frequency at the closing store
\end{itemize}
 
Results
\begin{itemize}
    \item Highly exposed households more likely to purchase something \textbf{other than their usual brand}
    \item Over time, new habits form for a new brand
    \item Effect driven by the most-exposed group (E4)
\end{itemize}
 
Mechanism
\begin{itemize}
    \item Closure removes contextual cues $\Rightarrow$ more \textbf{deliberative} decision-making $\Rightarrow$ exploration of previously ignored options
\end{itemize}
\end{frame}
 
 
% ------ Paper 4: "Store Closed" MSI WP ------
\begin{frame}{\textbf{``Store Closed''} (\textit{WP, 2023})}
\footnotesize
Qiaoni Shi, J. Jeffrey Inman, Dinesh Gauri
\vspace{4pt}
 
Setting \& Data
\begin{itemize}
    \item Two closure events: (1) mass merchandiser closing 100+ U.S.\ stores in 2016; (2) warehouse club closing 50+ stores in 2018
    \item Nielsen household panel: all trips, quantities, prices, store identities; 609 treated households
\end{itemize}
 
Identification
\begin{itemize}
    \item DiD with 20 months of data (10 pre, 10 post)
    \item Key moderator: distance to the nearest remaining open store of the same chain
\end{itemize}
 
Results
\begin{itemize}
    \item Most households \textbf{stayed with the retailer}
    \item Consumers bore closure costs via \textbf{greater travel} rather than switching
    \item \textbf{Distance} strongly moderates: farther households are more likely to switch retailers
    \item Gap: focuses on \textbf{where} customers shop, not \textbf{what} they buy
\end{itemize}
\end{frame}
 
 
% ------ Paper 3: Kotschedoff et al. (2025) ------
\begin{frame}{\textbf{Persistence of Grocery Shopping Behavior} (\textit{QME, 2025})}
\footnotesize
Marco J. W. Kotschedoff, Liliana Kowalczyk, and Els Breugelmans
\vspace{4pt}
 
Setting \& Data
\begin{itemize}
    \item European grocery chain: \textbf{one-week strike} $\Rightarrow$ temporary closure of a large fraction of stores
    \item Household purchase panel data (spending, store visits, category-level)
\end{itemize}
 
Identification
\begin{itemize}
    \item DiD: treated households vs.\ control households
    \item Calibrated structural simulations to back out implied state dependence
\end{itemize}
 
Results
\begin{itemize}
    \item During closure: consumers shift spending to competing chains
    \item After reopening: consumers \textbf{quickly revert} to pre-strike behavior $\Rightarrow$ consistent with \textbf{low state dependence}
    \item Perishable products: brief spending dip post-reopening; non-perishable: no change
\end{itemize}

\end{frame}
 
\end{comment}

\end{document}
